\chapter{2004年天津卷（文）}
\section{单选题}
\begin{problem}
设集合 \(P = \{1, 2, 3, 4, 5, 6\}\)，\(Q = \{x \in \R \mid 2 \leqslant x \leqslant 6\}\)，那么下列结论正确的是（\quad）

\choices
    {\(P\cap Q = P\)}
    {\(P\cap Q\supseteq Q\)}
    {\(P\cup Q = Q\)}
    {\(P\cap Q\nsupseteq P\)}
\answer{D}
\begin{solution}
由 \(Q=\left[2,6\right]\)，得
\(P\cap Q=\{2,3,4,5,6\}\)，这是 \(P\) 的真子集，所以 \(P\cap Q\ne P\)，且 \(P\cap Q\nsupseteq Q\). 又因为 \(1\in P\) 且 \(1\notin Q\)，所以 \(P\cup Q\ne Q\)，并且 \(P\cap Q\nsupseteq P\). 因此正确结论为 \(P\cap Q\nsupseteq P\)，与答案一致.
\end{solution}
\end{problem}

\begin{problem}
不等式 \(\frac{x - 1}{x} \geqslant 2\) 的解集为（\quad）

\choices
    {\([-1,0)\)}
    {\([-1, +\infty)\)}
    {\((-\infty, -1]\)}
    {\((-\infty, -1] \cup (0, +\infty)\)}
\answer{A}
\begin{solution}
分母不能为零，所以 \(x\ne0\). 移项并通分，原不等式等价于
\(\dfrac{x-1}{x}-2=\dfrac{-x-1}{x}\geqslant0\)，即 \(\dfrac{x+1}{x}\leqslant0\). 临界点为 \(-1\) 和 \(0\)，作符号判断可知在 \([-1,0)\) 上成立，在其余区间不成立，其中 \(-1\) 取等号而 \(0\) 不在定义域内. 因此解集为 \([-1,0)\)，故选 A.
\end{solution}
\end{problem}

\begin{problem}
对任意实数 \(a\)，\(b\)，\(c\) 在下列命题中，真命题是（\quad）

\choices
    {“\(ac > bc\)”是“\(a > b\)”的必要条件}
    {“\(ac = bc\)”是“\(a = b\)”的必要条件}
    {“\(ac > bc\)”是“\(a > b\)”的充分条件}
    {“\(ac = bc\)”是“\(a = b\)”的充分条件}
\answer{B}
\begin{solution}
因为 \(ac>bc\) 等价于 \(c(a-b)>0\)，当 \(a>b\) 且 \(c=0\) 时，\(ac=bc\)，所以“\(ac>bc\)”不是“\(a>b\)”的必要条件；当 \(c<0\) 时，\(ac>bc\) 反而推出 \(a<b\)，所以它也不是“\(a>b\)”的充分条件.

若 \(a=b\)，则无论 \(c\) 取何值都有 \(ac=bc\)，所以“\(ac=bc\)”是“\(a=b\)”的必要条件. 反过来取 \(c=0\) 且 \(a\ne b\)，仍有 \(ac=bc\)，故它不是充分条件. 因此正确命题是“\(ac=bc\)”是“\(a=b\)”的必要条件，与答案一致.
\end{solution}
\end{problem}

\begin{problem}
若平面向量 \(\bs{b}\) 与向量 \(\bs{a} = (1, -2)\) 的夹角是 \(180^{\circ}\)，且 \(\left|\bs{b}\right| = 3\sqrt{5}\)，则 \(\bs{b} =\)（\quad）

\choices
    {\((-3,6)\)}
    {\((3, - 6)\)}
    {\((6, - 3)\)}
    {\((-6,3)\)}
\answer{A}
\begin{solution}
因为两个向量的夹角为 \(180^{\circ}\)，所以 \(\bs{b}\) 与 \(\bs{a}\) 方向相反，设 \(\bs{b}=t\bs{a}\)，其中 \(t<0\). 又 \(\left|\bs{a}\right|=\sqrt{1^2+(-2)^2}=\sqrt5\)，故
\(\left|t\right|\sqrt5=3\sqrt5\)，从而 \(\left|t\right|=3\). 结合 \(t<0\)，得 \(t=-3\)，所以
\(\bs{b}=-3(1,-2)=(-3,6)\). 故选 A.
\end{solution}
\end{problem}

\begin{problem}
设 \(P\) 是双曲线 \(\frac{x^2}{a^2} - \frac{y^2}{9} = 1\) 上一点，双曲线的一条渐近线方程为 \(3x - 2y = 0\)，\(F_1\)，\(F_2\) 分别是双曲线的左，右焦点，若 \(|PF_1| = 3\)，则 \(|PF_2| =\)（\quad）

\choices
    {\(1\) 或 \(5\)}
    {\(6\)}
    {\(7\)}
    {\(9\)}
\answer{C}
\begin{solution}
双曲线的标准形式中，半实轴为 \(a\)，半虚轴为 \(b\)，渐近线为 \(y=\pm\dfrac{b}{a}x\). 题中分母 \(y^2\) 的系数给出 \(b=3\)，而渐近线 \(3x-2y=0\) 即 \(y=\dfrac32x\)，所以 \(\dfrac{b}{a}=\dfrac32\)，得 \(a=2\).

对双曲线上任一点 \(P\)，有 \(\left|PF_1-PF_2\right|=2a=4\). 已知 \(PF_1=3\)，故 \(\left|3-PF_2\right|=4\). 距离必须为正，因此只能取 \(PF_2=7\)，故选 C.
\end{solution}
\end{problem}

\begin{problem}
若函数 \( f(x) = \log_a x \) （\( 0 < a < 1 \)）在区间 \([a, 2a]\) 上的最大值是最小的 \(3\) 倍，则 \( a = \)（\quad）

\choices
    {\(\frac{\sqrt{2}}{4}\)}
    {\(\frac{\sqrt{2}}{2}\)}
    {\(\frac{1}{4}\)}
    {\(\frac{1}{2}\)}
\answer{A}
\begin{solution}
因为 \(0<a<1\)，函数 \(f(x)=\log_a x\) 在 \((0,+\infty)\) 上单调递减，所以在 \([a,2a]\) 上的最大值和最小值分别为
\(f(a)=1\) 和 \(f(2a)=\log_a2+1\). 由题意，
\(1=3\left(\log_a2+1\right)\)，所以 \(\log_a2=-\dfrac23\). 于是
\(a^{-\frac23}=2\)，即 \(a^{\frac23}=\dfrac12\)，从而
\(a=\left(\dfrac12\right)^{\frac32}=\dfrac{\sqrt2}{4}\). 故选 A.
\end{solution}
\end{problem}

\begin{problem}
若过定点 \(M(-1,0)\) 且斜率为 \(k\) 的直线与圆 \(x^{2} + 4x + y^{2} - 5 = 0\) 在第一象限内的部分有交点，则 \(k\) 的取值范围是（\quad）

\choices
    {\(0 < k < \sqrt{5}\)}
    {\(-\sqrt{5} < k < 0\)}
    {\(0 < k < \sqrt{13}\)}
    {\(0 < k < 5\)}
\answer{A}
\begin{solution}
圆的方程化为 \((x+2)^2+y^2=9\)，圆心为 \((-2,0)\)，半径为 \(3\). 过 \(M(-1,0)\) 且斜率为 \(k\) 的直线为 \(y=k(x+1)\). 要与第一象限内的圆弧相交，必须有 \(k>0\)，因为交点需满足 \(x>0\) 和 \(y>0\).

将直线代入圆的方程，得
\((x+2)^2+k^2(x+1)^2=9\). 在 \(x\geqslant0\) 时，左端随 \(x\) 增大而增大；当 \(x=0\) 时左端为 \(4+k^2\)，当 \(x=1\) 时左端为 \(9+4k^2>9\). 因此存在 \(0<x<1\) 的交点当且仅当 \(4+k^2<9\)，即 \(k<\sqrt5\). 端点 \(k=0\) 时交点在 \(x\) 轴上，端点 \(k=\sqrt5\) 时交点为 \((0,\sqrt5)\)，均不在第一象限内. 因此 \(0<k<\sqrt5\)，故选 A.
\end{solution}
\end{problem}

\begin{problem}
如图，定点 \(A\) 和 \(B\) 都在平面 \(\alpha\) 内，定点 \(P\notin \alpha\)，\(PB \perp \alpha\)，\(C\) 是 \(\alpha\) 内异于 \(A\) 和 \(B\) 的动点，且 \(PC \perp AC\). 那么，动点 \(C\) 在平面 \(\alpha\) 内的轨迹是（\quad）

\bitmapfigure[max width=0.36\linewidth]{img/2004/tianjin_liberal/q8_fig1.png}

\choices
    {一条线段，但要去掉两个点}
    {一个圆，但要去掉两个点}
    {一个椭圆，但要去掉两个点}
    {半圆，但要去掉两个点}
\answer{B}
\begin{solution}
在平面 \(\alpha\) 内取 \(B\) 为原点，令 \(\overrightarrow{BA}=\bs{a}\)，\(\overrightarrow{BC}=\bs{c}\)，并设 \(PB=h>0\). 由于 \(PB\perp\alpha\)，向量 \(\overrightarrow{CP}\) 的平面内分量为 \(-\bs{c}\)，而 \(\overrightarrow{CA}=\bs{a}-\bs{c}\). 条件 \(PC\perp AC\) 给出
\(\overrightarrow{CP}\cdot\overrightarrow{CA}=(-\bs{c})\cdot(\bs{a}-\bs{c})=0\)，即 \(\bs{c}\cdot(\bs{c}-\bs{a})=0\).

配方得
\(\left|\bs{c}-\dfrac12\bs{a}\right|^2=\dfrac14\left|\bs{a}\right|^2\). 因此 \(C\) 在以 \(AB\) 为直径的圆上. 题设规定 \(C\) 异于 \(A\) 和 \(B\)，所以要从该圆上去掉这两个点，故选 B.
\end{solution}
\end{problem}

\begin{problem}
函数 \( y = 3^{x^2 - 1} \)（\(-1 \leqslant x < 0\)）的反函数是（\quad）

\choices
    {\( y = \sqrt{1 + \log_{3} x} \)（\(x \geqslant \frac{1}{3}\)）}
    {\( y = -\sqrt{1 + \log_{3} x} \)（\(x \geqslant \frac{1}{3}\)）}
    {\( y = \sqrt{1 + \log_{3} x} \)（\(\frac{1}{3} < x \leqslant 1\)）}
    {\( y = -\sqrt{1 + \log_{3} x} \)（\(\frac{1}{3} < x \leqslant 1\)）}
\answer{D}
\begin{solution}
由 \(-1\leqslant x<0\)，可知 \(0<x^2\leqslant1\)，所以
\(-1<x^2-1\leqslant0\)，从而函数值满足 \(\dfrac13<y\leqslant1\). 由
\(y=3^{x^2-1}\) 得 \(\log_3y=x^2-1\)，即 \(x^2=1+\log_3y\). 原函数的定义域在负半轴上，故反函数应取负平方根：
\(y=-\sqrt{1+\log_3x}\)，且 \(\dfrac13<x\leqslant1\). 故选 D.
\end{solution}
\end{problem}

\begin{problem}
函数 \( y = 2\sin \left(\frac{\pi}{6} - 2x\right) \)（\(x \in [0, \pi]\)）为增函数的区间是（\quad）

\choices
    {\(\left[0, \frac{\pi}{3}\right]\)}
    {\(\left[\frac{\pi}{12}, \frac{7\pi}{12}\right]\)}
    {\(\left[\frac{\pi}{3}, \frac{5\pi}{6}\right]\)}
    {\(\left[\frac{5\pi}{6}, \pi\right]\)}
\answer{C}
\begin{solution}
求导得 \(y'=-4\cos\left(\dfrac\pi6-2x\right)\). 函数递增时需满足
\(\cos\left(\dfrac\pi6-2x\right)<0\). 在 \(x\in[0,\pi]\) 内解得
\(\dfrac\pi3<x<\dfrac{5\pi}{6}\). 在两个端点处导数为零，且导数在区间内部为正，所以函数在 \(\left[\dfrac\pi3,\dfrac{5\pi}{6}\right]\) 上为增函数，故选 C.
\end{solution}
\end{problem}

\begin{problem}
如图，在长方体 \(ABCD - A_{1}B_{1}C_{1}D_{1}\) 中，\(AB = 6\)，\(AD = 4\)，\(AA_{1} = 3\). 分别过 \(BC\)，\(A_{1}D_{1}\) 的两个平行截面将长方体分成三部分，其体积分别记为 \(V_{1} = V_{AEA_{1}-DFD_{1}}\)，\(V_{2} = V_{EBE_{1}A_{1}-FCF_{1}D_{1}}\)，\(V_{3} = V_{B_{1}E_{1}B-C_{1}F_{1}C}\). 若 \(V_{1}:V_{2}:V_{3} = 1:4:1\)，则截面 \(A_{1}EFD_{1}\) 的面积为（\quad）

\bitmapfigure[max width=0.42\linewidth]{img/2004/tianjin_liberal/q11_fig1.png}

\choices
    {\(4\sqrt{10}\)}
    {\(8\sqrt{3}\)}
    {\(4\sqrt{13}\)}
    {\(16\)}
\answer{C}
\begin{solution}
在正视图 \(ABB_1A_1\) 中，设 \(AE=x\). 两个截面平行，所以它们在正视图中的截线 \(A_1E\) 与 \(BE_1\) 平行；由两侧直角三角形的高都为 \(AA_1=BB_1=3\)，可得 \(B_1E_1=AE=x\). 同理，在后视图中 \(DF=x\)，且 \(C_1F_1=x\).

沿 \(AD\) 方向观察，三个部分的截面面积分别为两个全等直角三角形和中间剩余部分：
\(S_1=\dfrac12\cdot3x=\dfrac{3x}{2}\)，
\(S_3=\dfrac{3x}{2}\)，\(S_2=6\cdot3-S_1-S_3=18-3x\). 三部分沿同一方向延伸长度 \(AD=4\)，所以体积比等于截面面积比. 由题意
\(\dfrac{3x}{2}:\left(18-3x\right):\dfrac{3x}{2}=1:4:1\)，得
\(18-3x=4\cdot\dfrac{3x}{2}=6x\)，所以 \(x=2\).

截面 \(A_1EFD_1\) 是以 \(AD=4\) 和 \(A_1E\) 为两条互相垂直的边的矩形，且
\(A_1E=\sqrt{AA_1^2+AE^2}=\sqrt{3^2+2^2}=\sqrt{13}\). 因此其面积为 \(4\sqrt{13}\)，故选 C.
\end{solution}
\end{problem}

\begin{problem}
定义在 \(\R\) 上的函数 \(f(x)\) 既是偶函数又是周期函数，若 \(f(x)\) 的最小正周期是 \(\pi\)，且当 \(x \in \left[0, \frac{\pi}{2}\right]\) 时，\(f(x) = \sin x\)，则 \(f\left(\frac{5\pi}{3}\right)\) 的值为（\quad）

\choices
    {\(-\frac{1}{2}\)}
    {\(\frac{1}{2}\)}
    {\(-\frac{\sqrt{3}}{2}\)}
    {\(\frac{\sqrt{3}}{2}\)}
\answer{D}
\begin{solution}
由周期为 \(\pi\)，得
\(f\left(\dfrac{5\pi}{3}\right)=f\left(\dfrac{2\pi}{3}\right)\). 又因为 \(f\) 是偶函数，
\(f\left(\dfrac{2\pi}{3}\right)=f\left(-\dfrac{2\pi}{3}\right)=f\left(\dfrac\pi3\right)\)，其中最后一次等号利用了周期 \(\pi\). 由于 \(\dfrac\pi3\in\left[0,\dfrac\pi2\right]\)，所以
\(f\left(\dfrac{5\pi}{3}\right)=f\left(\dfrac\pi3\right)=\sin\dfrac\pi3=\dfrac{\sqrt3}{2}\). 故选 D.
\end{solution}
\end{problem}

\section{填空题}
\begin{problem}
某工厂生产 \(A\)，\(B\)，\(C\) 三种不同型号的产品，产品数量之比依次为 \(2:3:5\)，现用分层抽样方法抽出一个容量为 \(n\) 的样本，样本中 \(A\) 种型号产品有 \(16\) 件. 那么此样本的容量 \(n\) =\fillinblank{}.
\answer{80}
\begin{solution}
分层抽样时样本中各型号产品的数量比等于总体数量比. 设样本中三种产品的数量分别为 \(2t\)，\(3t\)，\(5t\)，则 \(2t=16\)，得 \(t=8\). 因此样本容量为
\(n=2t+3t+5t=10t=80\).
\end{solution}
\end{problem}

\begin{problem}
已知向量 \(\bs{a}=(1,1)\)，\(\bs{b}=(2,-3)\)，若 \(k\bs{a}-2\bs{b}\) 与 \(\bs{a}\) 垂直，则实数 \(k\) 等于\fillinblank{}.
\answer{\(-1\)}
\begin{solution}
由 \(k\bs{a}-2\bs{b}=(k,k)-(4,-6)=(k-4,k+6)\). 垂直条件给出
\((k-4,k+6)\cdot(1,1)=0\)，即 \(k-4+k+6=0\). 所以 \(2k+2=0\)，解得 \(k=-1\).
\end{solution}
\end{problem}

\begin{problem}
如果过两点 \(A(a,0)\) 和 \(B(0,a)\) 的直线与抛物线 \(y = x^{2} - 2x - 3\) 没有交点，那么实数 \(a\) 的取值范围是\fillinblank{}.
\answer{\(\left(-\infty,-\frac{13}{4}\right)\)}
\begin{solution}
过 \(A(a,0)\)，\(B(0,a)\) 的直线方程为 \(x+y=a\)，即 \(y=a-x\). 将它与抛物线方程联立，得
\(a-x=x^2-2x-3\)，整理为 \(x^2-x-a-3=0\). 直线与抛物线没有交点，等价于这个关于 \(x\) 的方程没有实根，因此判别式满足
\(\Delta=(-1)^2-4(-a-3)=4a+13<0\). 所以 \(a< -\dfrac{13}{4}\)，即 \(a\in\left(-\infty,-\dfrac{13}{4}\right)\).
\end{solution}
\end{problem}

\begin{problem}
从 \(0\)，\(1\)，\(2\)，\(3\)，\(4\)，\(5\) 中任取 \(3\) 个数字，组成没有重复数字的三位数，其中能被 \(5\) 整除的三位数共有\fillinblank{} 个.（用数字作答）
\answer{36}
\begin{solution}
能被 \(5\) 整除的三位数个位必须是 \(0\) 或 \(5\).

当个位为 \(0\) 时，百位从 \(1\)，\(2\)，\(3\)，\(4\)，\(5\) 中选取，有 \(5\) 种；十位从剩下的 \(4\) 个数字中选取，有 \(4\) 种，共有 \(5\cdot4=20\) 个.

当个位为 \(5\) 时，百位不能为 \(0\) 或 \(5\)，有 \(4\) 种；十位从剩下的 \(4\) 个数字中选取，共有 \(4\cdot4=16\) 个. 因此共有 \(20+16=36\) 个.
\end{solution}
\end{problem}

\section{解答题}
\begin{problem}
已知 \(\tan \left(\frac{\pi}{4} + \alpha\right) = \frac{1}{2}\).

\begin{enumerate}
\item 求 \(\tan\alpha\) 的值；
\item 求 \(\frac{\sin 2\alpha - \cos^{2}\alpha}{1 + \cos 2\alpha}\) 的值.
\end{enumerate}
\end{problem}

\begin{solution}
\begin{enumerate}
\item 令 \(t=\tan\alpha\). 由正切的和角公式，得
\(\dfrac{1+t}{1-t}=\dfrac12\)，所以 \(2+2t=1-t\)，从而 \(\tan\alpha=t=-\dfrac13\).
\item 由 \(\tan\alpha=-\dfrac13\)，得
\(\sin2\alpha=\dfrac{2\tan\alpha}{1+\tan^2\alpha}=-\dfrac35\)，
\(\cos^2\alpha=\dfrac{1}{1+\tan^2\alpha}=\dfrac9{10}\)，
\(\cos2\alpha=\dfrac{1-\tan^2\alpha}{1+\tan^2\alpha}=\dfrac45\). 因此
\(1+\cos2\alpha=\dfrac95\ne0\)，原式等于
\(\dfrac{-\frac35-\frac9{10}}{\frac95}=-\dfrac56\).
\end{enumerate}
\end{solution}

\begin{problem}
从 \(4\) 名男生和 \(2\) 名女生中任选 \(3\) 人参加演讲比赛.

\begin{enumerate}
\item 求所选 \(3\) 人都是男生的概率；
\item 求所选 \(3\) 人中恰有 \(1\) 名女生的概率；
\item 求所选 \(3\) 人中至少有 \(1\) 名女生的概率.
\end{enumerate}
\end{problem}

\begin{solution}
\begin{enumerate}
\item 从 \(6\) 人中任选 \(3\) 人，共有 \(\mathrm C_6^3=20\) 种等可能结果. 所选 \(3\) 人都是男生的结果有 \(\mathrm C_4^3=4\) 种，所以所求概率为
\(\dfrac{\mathrm C_4^3}{\mathrm C_6^3}=\dfrac15\).
\item 所选 \(3\) 人中恰有 \(1\) 名女生的结果有 \(\mathrm C_2^1\mathrm C_4^2=12\) 种，所以所求概率为
\(\dfrac{\mathrm C_2^1\mathrm C_4^2}{\mathrm C_6^3}=\dfrac35\).
\item 所选 \(3\) 人中至少有 \(1\) 名女生是“没有女生”的对立事件，因此所求概率为
\(1-\dfrac15=\dfrac45\).
\end{enumerate}
\end{solution}

\begin{problem}
如图，在四棱锥 \(P - ABCD\) 中，底面 \(ABCD\) 是正方形，侧棱 \(PD \perp\) 底面 \(ABCD\)，\(PD = DC\)，\(E\) 是 \(PC\) 的中点.

\begin{enumerate}
\item 证明：\(PA \parallel\) 平面 \(EDB\)；
\item 求 \(EB\) 与底面 \(ABCD\) 所成的角的正切值.
\end{enumerate}

\bitmapfigure[max width=0.46\linewidth]{img/2004/tianjin_liberal/q19_fig1.png}
\end{problem}

\begin{solution}
取正方形的边长为 \(a\)，则 \(PD=DC=a\). 建立空间直角坐标系，取
\(D(0,0,0)\)，\(A(a,0,0)\)，\(B(a,a,0)\)，\(C(0,a,0)\)，\(P(0,0,a)\). 因为 \(E\) 是 \(PC\) 的中点，所以 \(E\left(0,\dfrac a2,\dfrac a2\right)\).

\begin{enumerate}
\item 在平面 \(EDB\) 内取
\(\overrightarrow{DB}=(a,a,0)\) 和 \(\overrightarrow{DE}=\left(0,\dfrac a2,\dfrac a2\right)\). 向量 \(\bs{n}=(1,-1,1)\) 同时垂直于这两个向量，故可作为平面 \(EDB\) 的一个法向量.

又 \(\overrightarrow{PA}=(a,0,-a)\)，所以 \(\overrightarrow{PA}\cdot\bs{n}=a-a=0\)，即直线 \(PA\) 的方向向量平行于平面 \(EDB\). 平面 \(EDB\) 过原点 \(D\)，且 \(P\) 不在该平面内，因为 \(P\) 不满足平面方程 \(x-y+z=0\). 因此 \(PA\parallel\) 平面 \(EDB\).
\item 设点 \(E\) 在底面 \(ABCD\) 内的射影为 \(H\)，则 \(H\left(0,\dfrac a2,0\right)\). 于是
\(EH=\dfrac a2\)，且 \(HB=\sqrt{a^2+\left(\dfrac a2\right)^2}=\dfrac{a\sqrt5}{2}\). 设直线 \(EB\) 与底面 \(ABCD\) 所成的角为 \(\theta\)，则 \(\tan\theta=\dfrac{EH}{HB}=\dfrac1{\sqrt5}=\dfrac{\sqrt5}{5}\).
\end{enumerate}
\end{solution}

\begin{problem}
设 \(\{a_{n}\}\) 是一个公差为 \(d\)（\(d \neq 0\)）的等差数列，它的前 \(10\) 项和 \(S_{10} = 110\) 且 \(a_{1}\)，\(a_{2}\)，\(a_{4}\) 成等比数列.

\begin{enumerate}
\item 证明 \(a_1 = d\)；
\item 求公差 \(d\) 的值和数列 \(\{a_{n}\}\) 的通项公式.
\end{enumerate}
\end{problem}

\begin{solution}
\begin{enumerate}
\item 因为 \(a_1\)，\(a_2\)，\(a_4\) 成等比数列，所以 \(a_2^2=a_1a_4\). 又因为数列是等差数列，\(a_2=a_1+d\)，\(a_4=a_1+3d\)，故
\((a_1+d)^2=a_1(a_1+3d)\). 化简得 \(d^2=a_1d\). 由 \(d\ne0\)，可得 \(a_1=d\).
\item 等差数列前 \(10\) 项的和为
\(S_{10}=10a_1+\dfrac{10\cdot9}{2}d=10a_1+45d\). 代入 \(a_1=d\) 和 \(S_{10}=110\)，得 \(55d=110\)，所以 \(d=2\)，且 \(a_1=2\). 因此
\(a_n=a_1+(n-1)d=2+2(n-1)=2n\)，即数列的通项公式为 \(a_n=2n\).
\end{enumerate}
\end{solution}

\begin{problem}
已知函数 \( f(x) = ax^3 + cx + d \)（\(a \neq 0\)）是 \(\R\) 上的奇函数，当 \(x = 1\) 时 \(f(x)\) 取得极值 \(-2\).

\begin{enumerate}
\item 求 \( f(x) \) 的单调区间和极大值；
\item 证明对任意 \(x_{1}\)，\(x_{2} \in (-1, 1)\)，不等式 \(|f(x_{1}) - f(x_{2})| < 4\) 恒成立.
\end{enumerate}
\end{problem}

\begin{solution}
\begin{enumerate}
\item 因为 \(f(x)\) 是奇函数，所以 \(f(0)=0\)，从而 \(d=0\). 又因为 \(x=1\) 是函数的极值点，所以 \(f(1)=-2\)，且 \(f'(1)=0\). 于是
\(a+c=-2\)，\(3a+c=0\)，解得 \(a=1\)，\(c=-3\). 因此 \(f(x)=x^3-3x\)，且
\(f'(x)=3x^2-3=3(x-1)(x+1)\).

当 \(x\in(-\infty,-1)\cup(1,+\infty)\) 时，\(f'(x)>0\)；当 \(x\in(-1,1)\) 时，\(f'(x)<0\). 所以 \(f(x)\) 在 \((-\infty,-1)\) 和 \((1,+\infty)\) 上单调递增，在 \((-1,1)\) 上单调递减，且极大值为 \(f(-1)=2\).
\item 由上面的单调性，且 \(-1<x_i<1\)（\(i=1\)，\(2\)），得 \(-2<f(x_i)<2\). 因此
\(|f(x_1)-f(x_2)|<4\)，结论得证.
\end{enumerate}
\end{solution}

\begin{problem}
椭圆的中心是原点 \(O\)，它的短轴长为 \(2\sqrt{2}\). 相应于焦点 \(F(c,0)\) （\(c > 0\)）的准线 \(l\) 与 \(x\) 轴相交于点 \(A\)，\(|OF| = 2|FA|\). 过点 \(A\) 的直线与椭圆相交于 \(P\)，\(Q\) 两点.

\begin{enumerate}
\item 求椭圆的方程及离心率；
\item 若 \(\overrightarrow{OP} \cdot \overrightarrow{OQ} = 0\)，求直线 \(PQ\) 的方程.
\end{enumerate}
\end{problem}

\begin{solution}
\begin{enumerate}
\item 设椭圆的长半轴为 \(a\)，短半轴为 \(b\). 由短轴长为 \(2\sqrt2\)，得 \(b=\sqrt2\)，所以椭圆方程可设为 \(\dfrac{x^2}{a^2}+\dfrac{y^2}{2}=1\)，且 \(c^2=a^2-b^2=a^2-2\).

焦点 \(F(c,0)\) 对应的准线方程为 \(x=\dfrac{a^2}{c}\)，故 \(A\left(\dfrac{a^2}{c},0\right)\). 因为 \(a>c>0\)，所以
\(|FA|=\dfrac{a^2}{c}-c=\dfrac{a^2-c^2}{c}=\dfrac{b^2}{c}=\dfrac2c\). 由 \(|OF|=2|FA|\)，得 \(c=\dfrac4c\)，从而 \(c=2\)，\(a^2=b^2+c^2=6\). 因此椭圆方程为
\(\dfrac{x^2}{6}+\dfrac{y^2}{2}=1\)，离心率为 \(e=\dfrac ca=\dfrac{2}{\sqrt6}=\dfrac{\sqrt6}{3}\).
\item 由（1）知 \(A=(3,0)\). 直线 \(x=3\) 与椭圆无交点，故设直线 \(PQ\) 的方程为 \(y=k(x-3)\). 令 \(t=x-3\)，则交点的参数 \(t_1\)，\(t_2\) 是方程
\(\dfrac{(3+t)^2}{6}+\dfrac{k^2t^2}{2}=1\)，即 \((1+3k^2)t^2+6t+3=0\) 的两个根. 因此
\(t_1+t_2=-\dfrac6{1+3k^2}\)，\(t_1t_2=\dfrac3{1+3k^2}\).

交点为 \(P(3+t_1,kt_1)\)，\(Q(3+t_2,kt_2)\)，由 \(\overrightarrow{OP}\cdot\overrightarrow{OQ}=0\)，得
\(0=(3+t_1)(3+t_2)+k^2t_1t_2=9+3(t_1+t_2)+(1+k^2)t_1t_2\). 代入根的和与积，得
\(9-\dfrac{18}{1+3k^2}+\dfrac{3(1+k^2)}{1+3k^2}=0\)，所以 \(-6+30k^2=0\)，即 \(k^2=\dfrac15\).

故 \(k=\pm\dfrac1{\sqrt5}\)，直线 \(PQ\) 的方程为
\(y=\dfrac1{\sqrt5}(x-3)\) 或 \(y=-\dfrac1{\sqrt5}(x-3)\)，即 \(x-\sqrt5y-3=0\) 或 \(x+\sqrt5y-3=0\).
\end{enumerate}
\end{solution}
